% ============================================================
%  TEMA 3 — APARTADO 1: PLANIFICACIÓN DE RRHH O DE PLANTILLA
%
%  RECORDATORIO AL EQUIPO:
%  La planificación de RRHH (o planificación de plantilla) consiste en
%  anticipar las necesidades futuras de personal de la organización y
%  definir las acciones necesarias para satisfacerlas.
%
%  NIVELES DE PLANIFICACIÓN:
%    1. OPERATIVA (corto plazo, 1 año): cubrir necesidades inmediatas,
%       absentismo, bajas, picos de demanda.
%    2. TÁCTICA (medio plazo, 1-3 años): ajuste de plantilla, formación,
%       movilidad interna.
%    3. ESTRATÉGICA (largo plazo, +3 años): alineada con la estrategia
%       de negocio; transforma la estructura de la plantilla para los
%       retos futuros (digitalización, nuevos mercados, etc.).
%
%  PLANIFICACIÓN ESTRATÉGICA DE RRHH:
%    - Está integrada con el plan estratégico de la empresa.
%    - Considera escenarios futuros (tecnología, regulación, demografía).
%    - Incluye planificación de la sucesión de posiciones clave.
%    - Se mide con indicadores (KPIs de plantilla, turnover, etc.).
%
%  PARA ESTA SECCIÓN INDICAD:
%  - ¿Existe un proceso formal de planificación de plantilla en CaixaBank?
%  - ¿Es operativa, táctica o estratégica (o las tres)?
%  - ¿Está integrada con el plan estratégico del grupo?
%  - ¿Está la empresa satisfecha? ¿Qué mejoraría?
% ============================================================

\section{Tema 3: Planificación de recursos humanos}

\subsection{Existencia y naturaleza de la planificación de RRHH en CaixaBank}

% TODO: CaixaBank publica planes estratégicos trienales (p.ej. Plan Estratégico
% 2022-2024). Comprobad si ese plan incluye objetivos específicos de gestión
% de personas / plantilla. La sección de Sostenibilidad de la Memoria Anual
% también suele incluir datos de plantilla y objetivos de RRHH.

La planificación de recursos humanos es el proceso mediante el cual la
organización prevé sus necesidades futuras de personal y define las acciones
necesarias —contratación, formación, reducción, reconversión— para alinear
la dotación de personal con los objetivos estratégicos del negocio.

\subsubsection{¿Existe planificación formal de RRHH en CaixaBank?}

[Descripción a completar con la información obtenida en la entrevista
y en los documentos públicos de CaixaBank (Memoria Anual, Plan Estratégico)]

\subsubsection{¿Es estratégica la planificación de RRHH?}

% Para que la planificación sea estratégica debe:
%   a) Estar integrada con el plan estratégico de negocio.
%   b) Tener un horizonte temporal de al menos 3 años.
%   c) Considerar escenarios de futuro (no solo extrapolar el pasado).
%   d) Involucrar a la alta dirección, no solo al departamento de RRHH.
%   e) Contar con indicadores de seguimiento (KPIs de plantilla).

\begin{table}[H]
  \centering
  \small
  \caption{Características de la planificación estratégica de RRHH en CaixaBank}
  \begin{tabularx}{\textwidth}{p{0.45\textwidth} p{0.12\textwidth} X}
    \toprule
    \textbf{Característica}                          & \textbf{¿Se cumple?} & \textbf{Observaciones} \\
    \midrule
    Integración con plan estratégico de negocio      & [Sí/No/Parcial] & \\
    Horizonte temporal $\geq$ 3 años                 & [Sí/No/Parcial] & \\
    Consideración de escenarios de futuro            & [Sí/No/Parcial] & \\
    Involucración de la alta dirección               & [Sí/No/Parcial] & \\
    KPIs de seguimiento definidos                    & [Sí/No/Parcial] & \\
    Planificación de la sucesión                     & [Sí/No/Parcial] & \\
    \bottomrule
  \end{tabularx}
\end{table}

\subsubsection{Impacto de la fusión con Bankia en la planificación de plantilla}

% La fusión CaixaBank-Bankia (2021) fue la mayor fusión bancaria de la historia
% de España. Implicó un ajuste de plantilla significativo (ERE 2021).
% Analizad cómo afectó a la planificación de RRHH: ¿fue reactiva o proactiva?
% ¿Cómo se gestionó la integración de plantillas?

[Descripción a completar]

\begin{notaequipo}
  Preguntad en la entrevista:
  \begin{enumerate}
    \item ¿Existe un plan de plantilla formal? ¿Quién lo elabora y con qué periodicidad?
    \item ¿Está vinculado al plan estratégico del grupo?
    \item ¿Cómo se gestionó la planificación de plantilla tras la fusión con Bankia?
    \item ¿Tienen planificación de sucesión para posiciones clave?
    \item ¿Qué KPIs de plantilla monitorizan (ratio empleado/oficina, rotación,
          cobertura de vacantes, etc.)?
  \end{enumerate}
\end{notaequipo}
